% Quelle: 2. Übung Lösungen.pdf
% Art: Übung mit Lösungen
% Themen: ISBN10 Prüfziffer, Restklassenring, Einheitengruppe, Lemma von Bezout, Inverse in Z_m
\section{Übung 2: Wissen aus den Lösungen}
% Hinweis: Die Quelle ist ein Übungsblatt mit Lösungen. Die folgenden Wissensinhalte
% stehen in der Quelle innerhalb der Lösungen (Aufgabe 1 bzw. Aufgabe 5) und sind dort
% ebenfalls vollständig wiedergegeben.

\subsection{ISBN10}
\begin{bemerkung}[Prüfziffer ISBN10, aus Aufgabe 1]
% KORRIGIERT: stand "c_{10} = \sum ..." ohne Modul; die Gleichheit gilt nur modulo 11
\[
c_{10} \equiv \sum_{i=1}^{9} i \cdot c_i \pmod{11}, \qquad c_{10} \in \{0, 1, \dots, 10\}
\]
\end{bemerkung}

\subsection{Lemma von Bezout, Inverse in $\Zm{m}$}
\begin{bemerkung}[aus Aufgabe 5]
% KORRIGIERT: stand "Für teilerfremde Zahlen a, m" ohne Angabe, woher a und m stammen
Für $a \in \Z$ und $m \in \N$, $m \ge 2$, mit $\ggT(a, m) = 1$ besitzt die Diophantische Gleichung
\[
a \cdot x + m \cdot y = 1
\]
nach dem Lemma von Bezout Lösungen $(x, y)$ in ganzen Zahlen.
% KORRIGIERT: stand "x = a^{-1} in Z_m" (x ist eine ganze Zahl, gemeint sind Restklassen)
Das bedeutet \quad $a \cdot x \equiv 1 \pmod{m}$, also $[x] = [a]^{-1}$ in $\Zm{m}$.
\end{bemerkung}
